% =============================================================
\chapter{สรุปและข้อเสนอแนะ}
\label{ch:summary}
% =============================================================

\section{สรุปผลการดำเนินงาน}

ระบบ Disney Lorcana PlayLab Cloud ทำงานได้ครบตามขอบเขตในบทที่ 1 และมีหลักฐานการวัดผลบน AWS ดังนี้

\begin{enumerate}
    \item \textbf{ระยะที่ 2:} หน้าเว็บ ระบบสมาชิก จัดการเด็ค ห้องเล่นเรียลไทม์ และโครงสร้าง Multi-AZ (VPC, ALB, ASG)
    \item \textbf{Content Caching:} ลดขนาดการรับส่งไฟล์ข้อมูลการ์ด 81.43\% และเวลาโหลดซ้ำจาก 602 เหลือ 86 มิลลิวินาที
    \item \textbf{Dead Letter Queue:} ข้อความที่ประมวลผลล้มเหลวเข้า DLQ ภายใน 102 วินาที โดยไม่ค้างในคิวหลัก
    \item \textbf{SNS Fan-out:} ผลการแข่งขันปรากฏในประวัติและสถิติของทั้งสองผู้เล่นภายใน 5 วินาที พร้อมป้องกันการนับซ้ำ
    \item \textbf{ElastiCache Cache-Aside:} เวลาประมวลผลฝั่ง server ของตารางผู้นำลดจาก 41 เหลือ 1 มิลลิวินาที (ค่ามัธยฐาน) และมี fallback เมื่อแคชใช้ไม่ได้
    \item \textbf{Failover และ Auto Scaling:} แทนเครื่องที่ล่มได้ใน 2.8 ถึง 3.9 นาที และขยายจาก 2 เป็น 4 เครื่องเมื่อใช้ Node.js cluster
    \item \textbf{ตัวขับเคลื่อนกฎเกม:} รองรับ 16 action types, 14 คีย์เวิร์ด, Ability 10 ประเภท, Shift, Location, Song และ Item ผ่านชุดทดสอบอัตโนมัติ 128 รายการ
\end{enumerate}

ค่าใช้งานเดือนกันยายนทั้งเดือน 13.49 ดอลลาร์ก่อนหักเครดิต อยู่ภายในงบ 50 ดอลลาร์

\section{บทเรียนที่ได้รับ}

\begin{enumerate}
    \item \textbf{ค่าใช้จ่ายขึ้นกับทรัพยากรที่คิดรายชั่วโมง:} ALB คิดเป็น 57\% ของค่าใช้งานเดือนกันยายน การลด EC2 เป็น 0 ไม่หยุดค่า ALB
    \item \textbf{ออกแบบเผื่อหลายเครื่อง:} WebSocket ที่เก็บ state ใน memory ใช้ไม่ได้เมื่อมี 2 เครื่อง ต้องย้ายสถานะไป DynamoDB
    \item \textbf{เป้าหมาย scaling ต้องใช้ CPU ได้ถึงเป้า:} Node.js แบบ process เดียวใช้ CPU ได้ราวครึ่งเดียวของ t3.micro เป้า 60\% จึงไม่ถูกกระตุ้นจนกว่าจะใช้ cluster
    \item \textbf{สคริปต์ต้องตรวจผลหลังรันทุกครั้ง:} บั๊กชื่อ key ตาราง (การทดลองที่ 5) ผ่านการสร้างโดยไม่มี error แต่ทำให้ Lambda เขียนตารางไม่ได้ จึงควรมีขั้นตรวจ schema หลังสร้าง
    \item \textbf{เอกสารต้องตรวจกับระบบจริง:} รายงานระยะที่ 2 ระบุขนาด ASG และการรัน Docker ที่ไม่ตรงกับสคริปต์ที่ใช้ deploy จริง รายงานฉบับนี้แก้ให้ตรงและอ้างไฟล์หลักฐานทุกค่า
\end{enumerate}

\section{ข้อเสนอแนะ}

\begin{enumerate}
    \item ย้าย rate limiter จาก memory ไปใช้ DynamoDB พร้อม TTL เพื่อแก้ปัญหา OWASP A07
    \item เพิ่มเซิร์ฟเวอร์ตรวจสถานะเกม (authoritative game state) เพื่อป้องกันผลการแข่งขันที่ถูกแก้จาก client
    \item หากมีโดเมน ให้เพิ่ม HTTPS ผ่าน ACM และ ALB listener พอร์ต 443
    \item เพิ่มสัดส่วน Ability ที่รองรับเต็มรูปแบบจาก 16.16\% โดยเริ่มจากกลุ่ม trigger ที่พบมาก (activated 412 ข้อ, on\_play 596 ข้อ)
    \item หากใช้นอก Learner Lab ให้พิจารณา CloudFront และ ElastiCache Serverless
    \item เพิ่มการตรวจ schema ตารางหลังสร้างใน \texttt{ensure\_pipeline.ps1}
\end{enumerate}
